\documentclass[10pt,a4paper]{amsart}
\usepackage[margin=19mm]{geometry}
\usepackage{amsmath,amssymb,mathtools}
\usepackage[scr=rsfs,cal=euler]{mathalfa}
\usepackage{xfrac}
\usepackage[unicode,hidelinks,bookmarks=false]{hyperref}
\newcommand{\ie}{i.e.\ }
\newcommand{\cf}{cf.\ }
\newcommand{\resp}{respectively\ }
\newcommand{\Subsection}{\subsection}
\providecommand{\nobreakdash}{\nobreak\hbox{-}}
\setlength{\emergencystretch}{2em}
\multlinegap0pt
\usepackage{amssymb}
\usepackage{mathtools}
\usepackage{xcolor}
\usepackage{enumerate}
\usepackage{tikz}
\usepackage{float}
\newtheorem{proposition}{Proposition}
\newtheorem{theorem}{Theorem}
\newcommand{\E}{\mathcal{E}}
\DeclareMathOperator{\spa}{\textup{\textsf{span}}}
\title{Further results on modularity in evolution algebras}
\author{Manuel Ladra, Andr\'es P\'erez-Rodr\'iguez}
\date{Source: arXiv:2604.01111v1 (CC BY 4.0)}
\begin{document}
\maketitle

\noindent\emph{Source and licence.} Further results on modularity in evolution algebras. Version arXiv:2604.01111v1 (CC BY 4.0), distributed by arXiv under the Creative Commons Attribution 4.0 International licence, http://creativecommons.org/licenses/by/4.0/, as recorded in the arXiv licence metadata for this exact version and shown on the abstract page https://arxiv.org/abs/2604.01111v1. Full licence notice: \texttt{licenses/arxiv-2604.01111-CC-BY-4.0.txt}.\par\smallskip

\noindent\textit{Editorial scope.} Counted unit: the article's theorem th:char\_sup\_reg (source0.tex lines 536--538). The article's complete original proof of it is delivered with it (source0.tex lines 539--547, one \texttt{proof} environment of 9 lines). No mathematics is written, repaired or re-derived: the statement and the proof are the article's own lines, in the article's own notation, and no retained line is substituted or re-flowed.  Retained uncounted as the source-local context the proof needs: lines 386--388; lines 454--456.  Nothing was omitted from any retained span, so the package carries no omission record.  No retained line contains a citation, so the wrapper carries no bibliography and no bibliography entry.  No retained line contains a figure environment, an image-inclusion command or drawing code, so no image is delivered and the package declares no asset dependency.  Editorial formatting only, in addition: an AMS \texttt{amsart} preamble that restates the article's own declarations and macros used by the retained lines, namely the environments \texttt{proposition}, \texttt{theorem} on separate counters, together with the standard packages those lines need.  The retained environments print in this document as Proposition 1, Theorem 1 in order of appearance, whereas the article prints the same items with the numbering of its own full document.  The complete native source of the article as distributed in the arXiv e-print for this exact version is bundled unchanged as \texttt{sources/197638/source0.tex}.
\begin{proposition}\label{prop:super_triangular}
Let $\mathcal{E}$ be a regular evolution algebra over any field $\mathbb{K}$. Then, $\mathcal{E}$ is supersolvable if and only if there exists a natural basis such that the structure matrix is lower triangular with ones on the diagonal.
\end{proposition}

\begin{theorem}\label{th:char_mod}
Let $\mathcal{E}$ be a three-dimensional supersolvable regular evolution algebra over $\mathbb{C}$. Then, $\mathcal{E}$ is modular if and only if it is isomorphic to $\mathcal{E}_\mathbb{C}^{\mathrm{reg}}\big(0,\frac{1}{4},\frac{1}{4}\big)$.
\end{theorem}

\begin{theorem}\label{th:char_sup_reg}
Let $\mathcal{E}$ be a supersolvable regular evolution algebra over $\mathbb{C}$. Then, $\mathcal{E}$ is modular if and only if $n\leq2$ or $\mathcal{E}\cong\mathcal{E}_\mathbb{C}^{\mathrm{reg}}\big(0,\frac{1}{4},\frac{1}{4}\big)$.
\end{theorem}

\begin{proof}
Since every algebra of dimension at most two is modular, and the three-dimensional case has been characterised in Theorem~\ref{th:char_mod}, it suffices to consider the case $n>3$.
	
Let $\mathcal{E}$ be an $n$-dimensional supersolvable regular evolution algebra with natural basis $B=\{e_1,\dots,e_n\}$ and structure matrix $M_B(\mathcal{E})=(a_{ij})$ as in Proposition~\ref{prop:super_triangular}. If $\mathcal{E}$ is modular, then the subalgebra $\spa\{e_1,e_2,e_3\}$ would also be modular, since modularity is preserved under subalgebras. By Theorem~\ref{th:char_mod}, this forces that $a_{21}=0$ and $a_{31}=a_{32}=\frac{1}{4}$.
Now consider the quotient algebra $\mathcal{E}/\spa\{e_1\}$. Its structure matrix with respect to the natural basis $B_{\E/\spa\{e_1\}}=\{\overline{e_2},\dots,\overline{e_n}\}$ is obtained by deleting the first row and the first column of $M_B(\mathcal{E})$. In particular, the entry $(2,1)$ of this matrix equals $\tfrac14$. Consequently, Theorem~\ref{th:char_mod} implies that the three-dimensional subalgebra $\spa\{\overline{e_2},\overline{e_3},\overline{e_4}\}$ is not modular. Hence this quotient algebra is not modular.
Since modularity is preserved under both subalgebras and quotients, we obtain a contradiction. 

Therefore no supersolvable regular evolution algebra of dimension greater than three can be modular, and the result follows.
\end{proof}

\end{document}
